\documentclass[11pt]{article}
\usepackage[a4paper,margin=1in]{geometry}
\usepackage{amsmath,amssymb,amsthm,microtype}
\newtheorem{theorem}{Theorem}[section]
\newtheorem{corollary}[theorem]{Corollary}
\theoremstyle{remark}\newtheorem{remark}[theorem]{Remark}
\title{The Universal Logarithmic Singularity of the Endpoint Numerator}
\author{Bin Wang}\date{July 2026}
\begin{document}\maketitle
\begin{abstract}
The sharp deterministic coefficient law identifies the endpoint radius but
does not by itself identify the boundary singularity.  Using the exact odd
anchor and the exponentially decaying even trace remainder, we prove that
the endpoint-scaled logarithm is exactly $\log(1-w)+w$ plus a function
analytic on a strictly larger disk.  A repository-certified block constant
gives an explicit lower radius $1.0376199\ldots$ for that remainder.  This is
an all-order deterministic theorem, not convergence of a noisy spectrum.
\end{abstract}

\section{Normalized anchors}
Put $q_*=(r_H\lambda)^{-1}$, $\rho_*=q_*^{-1}$, and
$b_n=a_n\rho_*^n$.  The odd formula of RH-263 gives
\[
 b_n=(1+\lambda^{-n})^{-1}\qquad(n\ge3\text{ odd}).
\]
For $n=2m$, RH-268 writes
\[
 b_{2m}=2\lambda^{2m}\operatorname{tr}(T^m)
       +\frac{1-2\lambda^{-m}}{1-\lambda^{-2m}}.
\]
If $m=3k+j$, the certified trace estimate is
\[
 |\lambda^{2m}\operatorname{tr}(T^m)|
 \le C_j\delta^k,\qquad \delta<0.801254.
\]

\section{Singularity theorem}
\begin{theorem}
There is a function $H_{\rm reg}$ analytic for
\[
 |w|<r_{\rm reg},\qquad
 r_{\rm reg}\ge0.801254^{-1/6}=1.0376199142\ldots,
\]
such that
\[
 -\sum_{n\ge2}\frac{a_n\rho_*^n}{n}w^n
 =\log(1-w)+w+H_{\rm reg}(w).
\]
\end{theorem}
\begin{proof}
Write $b_n=1+r_n$.  At odd orders, $r_n=O(\lambda^{-n})$.
The explicit even endpoint fraction differs from one by
$O(\lambda^{-m})=O(\lambda^{-n/2})$.  The trace term is
$O(\delta^{n/6})$ after separating the three residue classes of $m$.
Consequently $\sum r_nw^n/n$ converges normally on every compact subdisk of
radius $\min\{\lambda,\sqrt\lambda,\delta^{-1/6}\}>1$.  Finally,
\[
 -\sum_{n\ge2}\frac{w^n}{n}=\log(1-w)+w,
\]
which proves the claim.
\end{proof}

\begin{corollary}
The logarithmic branch point at $w=1$ is the complete nonanalytic endpoint
part of the deterministic numerator logarithm.  In particular, subtracting
$\log(1-w)+w$ strictly enlarges the analytic disk.
\end{corollary}

\section{Scope}
The result uses exact all-order coefficient identities and a certified
geometric trace remainder; no finite-order fit is promoted.  It neither
constructs a positive or integer spectral cloud nor proves endpoint $H^2$
convergence of the actual mismatch.  Gates A--E remain false/open.
\end{document}
